\documentclass[12pt, a4paper]{article}

\usepackage[a4paper, top = 1 in, bottom = 1 in, left = 1 in, right = 1 in]{geometry}

\usepackage{amsmath, amssymb, amsfonts}
\allowdisplaybreaks[4]

\usepackage{enumerate}
\usepackage{multirow}
\usepackage{hhline}
\usepackage{array}
\usepackage{longtable}
\usepackage{graphicx}
\usepackage{tabularray}
\usepackage{undertilde}
\usepackage{xcolor}
\usepackage{dingbat}
\usepackage{fontawesome5}
\usepackage[colorlinks=true, linkcolor=blue, urlcolor=red]{hyperref}
\usepackage{tasks}
\usepackage{bbding}
\usepackage{twemojis}
% how to use bull's eye ----- \scalebox{2.0}{\twemoji{bullseye}}
\usepackage{customdice}
% how to put dice face ------ \dice{2}
\usepackage{fontspec}
\usepackage{fancyhdr}
\usepackage{lastpage}
\usepackage{mathrsfs}


\pagestyle{fancy}
\fancyhf{}
\fancyfoot[C]{Page \thepage\ of \pageref{LastPage}}
\renewcommand{\headrulewidth}{0pt}



\title{MSMS 401 : Bayesian Statistics and Computation}
\author{{\fontsize{25}{20} \benyorithafont Ananda Biswas}}
\date{}

\newfontface\gloriahallelujahfont{GLORIAHALLELUJAH-REGULAR.TTF}

\newfontface\benyorithafont{Benyoritha-G334D.otf}


\begin{document}

\maketitle

\section*{\faArrowAltCircleRight[regular] \textcolor{blue}{Problem}}

Let $X_i \overset{\text{iid}}{\sim} \text{Geometric}(\theta)$ on support  $\{0,1,2,\ldots\} \forall i = 1(1)n $ and \textit{a priori} $\theta \sim \mathscr{B}_1(\alpha,\beta)$, both $\alpha$ and $\beta$ are unknown. Obtain empirical Bayes estimate of $\alpha$ and $\beta$ by MLII method.

\section*{\faArrowAltCircleRight[regular] \textcolor{blue}{Solution}}

$$f(x \mid \theta) = \theta(1-\theta)^x, \qquad x=0,1,2,\ldots,\quad 0<\theta<1.$$

$$\pi(\theta; \alpha, \beta) = \dfrac{1}{\mathcal{B}(\alpha,\beta)} \cdot \theta^{\alpha-1}(1-\theta)^{\beta-1}, \qquad 0<\theta<1.$$

\vspace{0.5cm}

Then the joint marginal likelihood will be

\begin{align*}
L(\alpha, \beta) &= \displaystyle{\int \limits_0^1 \prod \limits_{i=1}^n f(x_i \mid \theta) \cdot \pi(\theta; \alpha, \beta)} \, d\theta \\[0.5em]
&= \displaystyle{\int \limits_0^1 \prod \limits_{i=1}^n \theta(1-\theta)^{x_i} \cdot \dfrac{1}{\mathcal{B}(\alpha,\beta)} \cdot \theta^{\alpha-1}(1-\theta)^{\beta-1}} \, d\theta \\[0.5em]
&= \displaystyle{\int \limits_0^1 \theta^n (1-\theta)^{\sum \limits_{i=1}^n x_i} \cdot \dfrac{1}{\mathcal{B}(\alpha,\beta)} \cdot \theta^{\alpha-1}(1-\theta)^{\beta-1}} \, d\theta \\[0.5em]
&= \dfrac{1}{\mathcal{B}(\alpha,\beta)} \cdot \displaystyle{\int \limits_0^1 \theta^{\alpha + n - 1} (1 - \theta)^{\beta + \sum \limits_{i=1}^n x_i - 1}} \, d\theta \\[0.5em]
&= \dfrac{1}{\mathcal{B}(\alpha,\beta)} \cdot \mathcal{B}\left(\alpha + n,\beta + \sum \limits_{i=1}^n x_i \right) \\[0.7em]
&= \dfrac{\Gamma(\alpha+\beta)}{\Gamma(\alpha)\Gamma(\beta)} \cdot \dfrac{\Gamma(\alpha + n)\Gamma\left(\beta + \sum \limits_{i=1}^n x_i\right)}{\Gamma\left(\alpha + \beta + \sum \limits_{i=1}^n x_i + n\right)} \\[0.7em]
&= \dfrac{\Gamma(\alpha+\beta) \cdot \Gamma(\alpha+n) \cdot \Gamma\left(\beta + \sum \limits_{i=1}^n x_i\right)}{\Gamma(\alpha) \cdot \Gamma(\beta) \cdot \Gamma\left(\alpha + \beta + \sum \limits_{i=1}^n x_i + n\right)}
\end{align*}

\vspace{0.5cm}

The log-likelihood will be

\begin{align*}
\ell(\alpha, \beta) &= \ln \Gamma(\alpha+\beta) + \ln \Gamma(\alpha + n) + \ln \Gamma\left(\beta + \sum \limits_{i=1}^n x_i\right) \\
&\hspace{4cm} - \ln \Gamma(\alpha) - \ln \Gamma(\beta) - \ln \Gamma\left(\alpha + \beta + \sum \limits_{i=1}^n x_i + n\right).
\end{align*}

Here both $\alpha$ and $\beta$ are unknown. Partial differentiation of $\ell(\alpha, \beta)$ w.r.t. $\alpha$ and $\beta$ introduces \textbf{digamma} function. Hence, analytical solution of the equations

\begin{align*}
\dfrac{\partial}{\partial \alpha} \ell(\alpha, \beta) &= 0 \\[0.5em]
\dfrac{\partial}{\partial \beta} \ell(\alpha, \beta) &= 0
\end{align*}

is not possible and we resort to numerical methods for approximate solutions.
\end{document}